\chapter{System Design}

\section{System Architecture}
The development of HomeSOC necessitated a highly robust, fault-tolerant, three-tiered architecture. The objective was to facilitate the seamless localized collection of telemetry, its subsequent transmission, ingestion, long-term storage, heuristic processing, and final frontend visualisation. This chapter delineates the top-down design methodology that guided the construction of the overall system infrastructure. 

The fundamental components of the project were logically compartmentalised into three distinct layers:
\begin{enumerate}
    \item \textbf{The Edge Layer (Host Agent)}: Deployed directly onto the target hardware intended for monitoring.
    \item \textbf{The Core Backend (Ingestion Server \& Detection Engine)}: The central repository and analytical brain hosted securely physically or via cloud.
    \item \textbf{The Presentation Layer (Dashboard Interface)}: The web-based graphical controller accessed by the end-user.
\end{enumerate}

\begin{figure}[h]
    \centering
    % Placeholder for Overall Architecture Diagram
    \fbox{
        \begin{minipage}{0.8\textwidth}
            \vspace{2cm}
            \begin{center}
                \textit{[PLACEHOLDER: Insert Structural UML/Architecture Block Diagram here]} \\
                \vspace{0.5cm}
                File: \texttt{architecture\_diagram.png} \\
                \textit{Show Agent $\rightarrow$ Cloud API $\rightarrow$ Database $\leftarrow$ Frontend Browser}
            \end{center}
            \vspace{2cm}
        \end{minipage}
    }
    \caption{The macro three-tier architectural topology of the HomeSOC platform.}
    \label{fig:macro_arch}
\end{figure}

\section{System Flow Interaction}
The system architecture was constructed with explicit attention to modular execution and stateless transitions. The following hierarchy rigorously identifies how each functional block interacts with volatile data streams:

\begin{enumerate}
    \item \textbf{Hardware Telemetry Acquisition}: The agent commences the initial localized collection of Operating System (OS) resource statistics (CPU throttling, RAM overcommitment, Disk I/O waits) and scans physical/virtual authentication log directories (e.g., `/var/log/auth.log`).
    \item \textbf{Ingestion \& RESTful API Handshake}: The extracted artifacts are bundled into a serialized JSON payload and transmitted securely over standard HTTP/HTTPS POST requests toward the Flask application's `/api/logs` listener.
    \item \textbf{Asynchronous Heuristic Analysis}: Upon receipt, the backend immediately intercepts the raw data array, diverting it into the rule-based detection engine (`detection.py`) before standard archival, dramatically increasing detection velocity.
    \item \textbf{Alert Graduation and SQL State Change}: Conversion of mathematically flagged anomalies into formal `alerts`. These objects are appended sequentially into the SQLite storage matrices with associated metadata (host origin, severity classification, and precise temporal timestamps).
    \item \textbf{Dashboard Triage and User Fulfillment}: The vanilla Javascript frontend executes a continuous asynchronous polling function (`setInterval`) against the `/api/alerts` endpoint, injecting the latest JSON response directly into the visible browser Window DOM for active monitoring and subsequent human resolution.
\end{enumerate}

\section{System Flowchart Specification (UML)}
A structural sequence diagram is fundamentally necessary for illustrating the strict cyclical nature of telemetry collection across the network boundary diagram. Figure 3-2 visualizes the transaction layer.

\begin{figure}[h]
    \centering
    \begin{tikzpicture}[node distance=2.5cm, auto,
            agent/.style={rectangle, draw, fill=blue!10, text width=6em, text centered, rounded corners, minimum height=3em},
            server/.style={rectangle, draw, fill=red!10, text width=6em, text centered, rounded corners, minimum height=3em},
            db/.style={cylinder, draw, shape border rotate=90, aspect=0.25, fill=gray!20, text width=4em, text centered, minimum height=3em},
            client/.style={rectangle, draw, fill=green!10, text width=6em, text centered, rounded corners, minimum height=3em}]
            
        \node [agent] (node1) {Endpoint Agent};
        \node [server, right=of node1, xshift=1cm] (flask) {Flask API Web Server};
        \node [db, right=of flask, xshift=0.5cm] (sqlite) {SQLite Database};
        \node [client, below=of flask] (browser) {User Browser};
        
        % Agent to Server
        \draw[->, thick, dashed] (node1.east) -- node[above] {1. POST /logs (JSON)} (flask.west);
        
        % Server to DB (Write)
        \draw[->, thick] (flask.east) -- node[above] {2. INSERT INTO} node[below] {(Heuristic Trigger)} (sqlite.west);
        
        % Browser to Server (Poll)
        \draw[->, thick] (browser.north) -- node[left] {3. GET /alerts} (flask.south);
        
        % Server to Browser (Response)
        \draw[->, thick, dashed] (flask.south) to[bend left=30] node[right] {4. JSON Array} (browser.east);

    \end{tikzpicture}
    \caption{Unified Modeling Language (UML) structural sequence detailing the full API request lifecycle.}
    \label{fig:uml_seq}
\end{figure}

\section{Database Schema Design Block}
To persistently track the operational history of the network without overwhelming processing overhead or incurring significant licensing costs, a standalone SQLite database engine was permanently integrated directly into the `app.py` wrapper directory structure. 

Crucially, to prevent data concurrency deadlocks during simultaneous logging and querying by multiple users—a notorious weakness in basic SQLite setups deployed across network file systems (NFS)—the Write-Ahead Logging (WAL) pragmas were manually enabled within the `db.py` driver abstraction layer. 

The core schema comprises four critical relational tables linked via explicit foreign keys, establishing strict referential integrity.

\subsection{Relational Entity Mappings}
Table 3-1 presents the schema definition for the primary `alerts` archive, highlighting data types and functional descriptions.

\begin{table}[h]
\centering
\small
\begin{tabularx}{\textwidth}{|l|c|X|}
\hline
\textbf{Field Name} & \textbf{Data Type} & \textbf{Functional Description} \\ \hline
id & INTEGER & Primary Key (Auto-increment format) for the alert incident. \\ \hline
host\_id & INTEGER & Relational Foreign Key structurally linking to the source host in the `hosts` table. \\ \hline
rule\_name & TEXT & The nominal descriptor of the triggered detection heuristic (e.g., ``High CPU Utilization''). \\ \hline
severity & TEXT & Categorical severity ranking constraint: `Low`, `Medium`, `High`. \\ \hline
category & TEXT & Attack vector classification (e.g., `metrics`, `auth`, `process\_list`). \\ \hline
raw\_payload & JSON TEXT & The direct, immutable evidentiary telemetry string extracted verbatim from the original system logs. \\ \hline
status & TEXT & Workflow lifecycle state string identifier: `new`, `handled`, `resolved`. \\ \hline
timestamp & DATETIME & Temporal integer auto-generated on insertion (CURRENT\_TIMESTAMP). \\ \hline
\end{tabularx}
\caption{Detailed database schema mapping for the primary SIEM `alerts` management table.}
\end{table}

\section{Detection Engine Heuristic Logic}
The central intelligence and cognitive processing weight of the SOC ultimately resides in the `detection.py` micro-service. The analytical layer shuns resource-intensive Machine Learning operations in favor of a deterministic, rigid threshold-based calculation script. 

The engine actively unmarshals the incoming JSON transmission and extracts numeric integers corresponding to pre-mapped key-value dictionary combinations.
For example, standard resource anomaly checks iterate through the incoming system payload matrix in real-time. 

\subsection{Pseudo-Code Algorithmic Implementation}
The foundational logic for capturing CPU resource starvation natively is defined algorithmically as follows:

\begin{verbatim}
FUNCTION EvaluateLocalPayload(log_entry, database_rules)
    alerts_queue = []
    
    FOR rule IN database_rules:
        condition_logic = rule.get('logic_syntax')
        operator = condition_logic.get('operator') // e.g., '>', '<', 'contains'
        threshold = condition_logic.get('value')
        
        extracted_metric = ExtractTargetData(log_entry, rule.field)
        
        IF operator == '>' AND Type(extracted_metric) == Integer:
            IF extracted_metric > threshold:
                alerts_queue.Append({
                    'Triggered_Rule': rule.name,
                    'Evidence_Data': log_entry,
                    'Host': log_entry.origin
                })
        
        IF operator == 'contains' AND Type(extracted_metric) == Array:
             IF threshold IN extracted_metric: // Example: 'nmap' IN processes
                  alerts_queue.Append(...)
                  
    RETURN alerts_queue
\end{verbatim}

By architecturally decoupling the Rule Definitions (stored statically as SQL rows) from the Logic Engine execution loop, system administrators receive the distinct strategic advantage of freely updating security parameters globally without necessitating physical modifications or operational downtime to the underlying Python execution environment. This strict topological isolation dramatically increases the long-term maintainability, patchability, and functional modularity of the HomeSOC system.

\begin{figure}[h]
    \centering
    % Placeholder for Sandbox View Screenshot
    \fbox{
        \begin{minipage}{0.8\textwidth}
            \vspace{2cm}
            \begin{center}
                \textit{[PLACEHOLDER: Insert Screenshot of SQL Sandbox Tab UI Dashboard]} \\
                \vspace{0.5cm}
                File: \texttt{sandbox\_execution.png}
            \end{center}
            \vspace{2cm}
        \end{minipage}
    }
    \caption{The SQL analytical sandbox allowing real-time administrative access into the raw relational schema.}
    \label{fig:sandbox_ui}
\end{figure}

\section{REST API Mapping Configuration}
A comprehensive API matrix defines the endpoints interacting with the client-side logic. To formalize backend-frontend separation, Table 3-2 details the defined operational HTTP routing paths.

\begin{table}[h]
\centering
\small
\begin{tabularx}{\textwidth}{|X|c|X|c|}
\hline
\textbf{Endpoint URI Route} & \textbf{Method} & \textbf{Functionality Protocol Description} & \textbf{Auth Needs} \\ \hline
\texttt{/api/heartbeat} & POST & Registers host nodes dynamically into the `hosts` schema table. & None (MVP) \\ \hline
\texttt{/api/logs} & POST & Core telemetry ingestion endpoint executing detection heuristics. & None (MVP) \\ \hline
\texttt{/api/alerts/active} & GET & Real-time fetch sequence feeding the DOM dashboard arrays. & Yes \\ \hline
\texttt{/api/alerts/history} & GET & Retrieves resolved, handled, or suppressed historical alerts. & Yes \\ \hline
\texttt{/api/query} & POST & Direct raw SQLite execution sandbox query endpoint. & Strict Admin \\ \hline
\end{tabularx}
\caption{The structured REST API web routing matrix configuring host and interface communication layers.}
\label{tab:api_maps}
\end{table}

\section{Expanded Theoretical Framework: Database Architecture and SQLite B-Trees}
SQLite ( "S-Q-L-ite",  "sequel-ite") is a free and open-source relational database engine written in the C programming language. It is not a standalone application; rather, it is a library that software developers embed in their applications. As such, it belongs to the family of embedded databases. According to its developers, SQLite is the most widely deployed database engine, as it is used by several of the top web browsers, operating systems, mobile phones, and other embedded systems.
Many programming languages have bindings to the SQLite library. It generally follows PostgreSQL syntax, but does not enforce type checking by default. This means that one can, for example, insert a string into a column defined as an integer. Although it is a lightweight embedded database, SQLite implements most of the SQL standard and the relational model, including transactions and ACID guarantees. However, it omits many features implemented by other databases, such as materialized views and complete support for triggers and ALTER TABLE statements.


== History ==
D. Richard Hipp designed SQLite in the spring of 2000 while working for General Dynamics on contract with the United States Navy. Hipp was designing software used for a damage-control system aboard guided-missile destroyers; the damage-control system originally used HP-UX with an Informix database back-end. SQLite began as a Tcl extension.
In August 2000, version 1.0 of SQLite was released, with storage based on gdbm (GNU Database Manager). In September 2001, SQLite 2.0 replaced gdbm with a custom B-tree implementation, adding transaction capability. In June 2004, SQLite 3.0 added internationalization, manifest typing, and other major improvements, partially funded by America Online. In 2011, Hipp announced his plans to add a NoSQL interface to SQLite, as well as announcing UnQL, a functional superset of SQL designed for document-oriented databases.
SQLite is one of four formats recommended for long-term storage of datasets approved for use by the Library of Congress.


== Design ==
SQLite was designed to allow the program to be operated without installing a database management system or requiring a database administrator. Unlike client–server database management systems, the SQLite engine has no standalone processes with which the application program communicates. Instead, a linker integrates the SQLite library—statically or dynamically—into an application program which uses SQLite's functionality through simple function calls, reducing latency in database operations; for simple queries with little concurrency, SQLite performance profits from avoiding the overhead of inter-process communication.
Due to the serverless design, SQLite applications require less configuration than client–server databases. SQLite is called zero-configuration because configuration tasks such as service management, startup scripts, and password- or GRANT-based access control are unnecessary. Access control is handled through the file-system permissions of the database file. Databases in client–server systems use file-system permissions that give access to the database files only to the daemon process, which handles its locks internally, allowing concurrent writes from several processes.
SQLite stores the entire database, consisting of definitions, tables, indices, and data, as a single cross-platform file, allowing several processes or threads to access the same database concurrently. It implements this simple design by locking the database file during writing. Write access may fail with an error code, or it can be retried until a configurable timeout expires. SQLite read operations can be multitasked, though due to the serverless design, writes can only be performed sequentially. This concurrent access restriction does not apply to temporary tables, and it is relaxed in version 3.7 as write-ahead logging (WAL) enables concurrent reads and writes. Since SQLite has to rely on file-system locks, it is not the preferred choice for write-intensive deployments.
SQLite uses PostgreSQL as a reference platform. "What would PostgreSQL do" is used to make sense of the SQL standard. One major deviation is that, with the exception of primary keys, SQLite does not enforce type checking; the type of a value is dynamic and not strictly constrained by the schema (although the schema will trigger a conversion when storing, if such a conversion is potentially reversible). SQLite strives to follow Postel's rule.


== Features ==
SQLite implements most of the SQL-92 standard for SQL, but lacks some features. For example, it only partially provides triggers and cannot write to views (however, it provides INSTEAD OF triggers that provide this functionality). Its support of ALTER TABLE statements is limited.
SQLite uses an unusual type system for an SQL-compatible DBMS: instead of assigning a type to a column as in most SQL database systems, types are assigned to individual values; in language terms it is dynamically typed. Moreover, it is weakly typed in some of the same ways that Perl is: one can insert a string into an integer column (although SQLite will try to convert the string to an integer first, if the column's preferred type is integer). This adds flexibility to columns, especially when bound to a dynamically typed scripting language. However, the technique is not portable to other SQL products. A common criticism is that SQLite's type system lacks the data integrity mechanism provided by statically typed columns, although it can be emulated with constraints like CHECK(typeof(x)='integer'). In 2021, support for static typing was added through STRICT tables, which enforce datatype constraints for columns.
Tables normally include a hidden rowid index column, which provides faster access. If a table includes an INTEGER PRIMARY KEY column, SQLite will typically optimize it by treating it as an alias for the rowid, causing the contents to be stored as a strictly typed 64-bit signed integer and changing its behavior to be somewhat like an auto-incrementing column. SQLite includes an option to create a table without a rowid column, which can save disk space and improve lookup speed. WITHOUT ROWID tables are required to have a primary key.
SQLite supports foreign key constraints, although they are disabled by default and must be manually enabled with a PRAGMA statement.
Stored procedures are not supported; this is an explicit choice by the developers to favor simplicity, as the typical use case of SQLite is to be embedded inside a host application that can define its own procedures around the database.
SQLite does not have full Unicode support by default for backwards compatibility and due to the size of the Unicode tables, which are larger than the SQLite library. Full support for Unicode case-conversions can be enabled through an optional extension.
SQLite supports full-text search through its FTS5 loadable extension, which allows users to efficiently search for a keyword in a large number of documents similar to how search engines search webpages.
SQLite includes support for working with JSON through its json1 extension, which is enabled by default since 2021. SQLite's JSON functions can handle JSON5 syntax since 2023. In 2024, SQLite added support for JSONB, a binary serialization of SQLite's internal representation of JSON. Using JSONB allows applications to avoid having to parse the JSON text each time it is processed and saves a small amount of disk space.
In May 2025, the 25th‑anniversary release SQLite 3.50.0 introduced additional features, including new Unicode functions (unistr() and unistr\_quote()), a new API (sqlite3\_setlk\_timeout()) for setting lock timeouts, improved command‑line tools and rsync utility enhancements, and optimized JSONB.
The maximum supported size for an SQLite database file is 281 terabytes.


== Development and distribution ==
SQLite's code is hosted with Fossil, a distributed version control system that uses SQLite as a local cache for its non-relational database format, and SQLite's SQL as an implementation language.
SQLite is public domain, but not "open-contribution", with the website stating "the project does not accept patches from people who have not submitted an affidavit dedicating their contribution into the public domain." Instead of a code of conduct, the founders have adopted a code of ethics based on the Rule of St. Benedict.
A standalone command-line shell program called sqlite3 is provided in SQLite's distribution. It can be used to create a database, define tables, insert and change rows, run queries and manage an SQLite database file. It also serves as an example for writing applications that use the SQLite library.
SQLite uses automated regression testing prior to each release. Over 2 million tests are run as part of a release's verification. The SQLite library has 156,000 lines of source code, while all the test suites combined add up to 92 million lines of test code. SQLite's tests simulate a number of exceptional scenarios, such as power loss and I/O errors, in addition to testing the library's functionality. Starting with the August 10, 2009 release of SQLite 3.6.17, SQLite releases have 100\% branch test coverage, one of the components of code coverage. SQLite has four different test harnesses: the original public-domain TCL tests, the proprietary C-language TH3 test suite, the SQL Logic Tests, which check SQLite against other SQL databases, and the dbsqlfuzz proprietary fuzzing engine.


== Notable uses ==


=== Operating systems ===
SQLite is included by default in:

Android
BlackBerry 10 OS
Fedora Linux where it is used by the rpm core package management system
FreeBSD where starting with 10-RELEASE version in January 2014, it is used by the core package management system.
illumos
iOS
Mac OS X 10.4 onwards (Apple adopted it as an option in macOS's Core Data API from the original implementation)
Maemo
MeeGo
MorphOS 3.10 onwards
NetBSD
NixOS where it is used by the Nix core package management system
Red Hat Enterprise Linux where it is used in the same way as Fedora, from which Red Hat Enterprise Linux is derived
Solaris 10 where the Service Management Facility database is serialized for booting.
Symbian OS
Tizen
webOS
Windows 10 onwards


=== Middleware ===
ADO.NET adapter, initially developed by Robert Simpson, is maintained jointly with the SQLite developers since April 2010.
ODBC driver has been developed and is maintained separately by Christian Werner. Werner's ODBC driver is the recommended connection method for accessing SQLite from OpenOffice.org.
COM (ActiveX) wrapper making SQLite accessible on Windows to scripted languages such as JScript and VBScript. This adds SQLite database capabilities to HTML Applications (HTA).


=== Web browsers ===
The browsers Google Chrome, Opera, Safari and the Android Browser all allow for storing information in, and retrieving it from, an SQLite database within the browser, using the official SQLite Wasm (WebAssembly) build, or using the Web SQL Database technology, although the latter is becoming deprecated (namely superseded by SQLite Wasm or by IndexedDB). Internally, these Chromium based browsers use SQLite databases for storing configuration data like site visit history, cookies, download history etc.
Mozilla Firefox and Mozilla Thunderbird store a variety of configuration data (bookmarks, cookies, contacts etc.) in internally managed SQLite databases. Until Firefox version 57 ("Firefox Quantum"), there was a third-party add-on that used the API supporting this functionality to provide a user interface for managing arbitrary SQLite databases.
Several third-party add-ons can make use of JavaScript APIs to manage SQLite databases.


=== Web application frameworks ===
Symfony
Laravel
Bugzilla
Django's default database management system
Drupal
Trac
Ruby on Rails's default database management system
web2py


=== Others ===
Adobe Systems uses SQLite as its file format in Adobe Lightroom, a standard database in Adobe AIR, and internally within Adobe Reader.
Apple Photos uses SQLite internally.
Audacity uses SQLite as its file format, as of version 3.0.0.
Evernote uses SQLite to store its local database repository in Windows.
Skype
WhatsApp
The Service Management Facility, used for service management within the Solaris and OpenSolaris operating systems
Dropbox client software
Intuit uses SQLite in QuickBooks and TurboTax
McAfee Antivirus
Flame (malware)
BMW iDrive satellite navigation system
TomTom GPS systems, for the NDS map data
Proxmox VE - Proxmox Cluster File System (pmxcfs)
Bentley Systems MicroStation
Bosch car multimedia systems
Airbus A350 flight system
Quicken Essentials and later versions of Quicken for Mac
Python standard library
Xojo IDE


== See also ==

Ordered key–value store
Comparison of relational database management systems
List of relational database management systems
MySQL
SpatiaLite


== Notes ==

A relational database (RDB) is a database  based on the relational model of data, as proposed by E. F. Codd in 1970.
A Relational Database Management System (RDBMS) is a type of database management system that stores data in a structured format using rows and columns.
Many relational database systems are equipped with the option of using SQL (Structured Query Language) for querying and updating the database.


== History ==
The concept of relational database was defined by E. F. Codd at IBM in 1970. Codd introduced the term relational in his research paper "A Relational Model of Data for Large Shared Data Banks". In this paper and later papers, he defined what he meant by relation. One well-known definition of what constitutes a relational database system is composed of Codd's 12 rules.
However, no commercial implementations of the relational model conform to all of Codd's rules, so the term has gradually come to describe a broader class of database systems, which at a minimum:

Present the data to the user as relations (a presentation in tabular form, i.e. as a collection of tables with each table consisting of a set of rows and columns);
Provide relational operators to manipulate the data in tabular form.
In 1974, IBM began developing System R, a research project to develop a prototype RDBMS.
The first system sold as an RDBMS was Multics Relational Data Store (June 1976). Oracle was released in 1979 by Relational Software, now Oracle Corporation. Ingres and IBM BS12 followed. Other examples of an RDBMS include IBM Db2, SAP Sybase ASE, and Informix. In 1984, the first RDBMS for Macintosh began being developed, code-named Silver Surfer, and was released in 1987 as 4th Dimension and known today as 4D.
The first systems that were relatively faithful implementations of the relational model were from:

University of Michigan – Micro DBMS (1969)
Massachusetts Institute of Technology (1971)
IBM UK Scientific Centre at Peterlee – IS1 (1970–72), and its successor, PRTV (1973–79).
The most common definition of an RDBMS is a product that presents a view of data as a collection of rows and columns, even if it is not based strictly upon relational theory. By this definition, RDBMS products typically implement some but not all of Codd's 12 rules.
A second school of thought argues that if a database does not implement all of Codd's rules (or the current understanding on the relational model, as expressed by Christopher J. Date, Hugh Darwen and others), it is not relational. This view, shared by many theorists and other strict adherents to Codd's principles, would disqualify most DBMSs as not relational. For clarification, they often refer to some RDBMSs as truly-relational database management systems (TRDBMS), naming others pseudo-relational database management systems (PRDBMS). 
As of 2009, most commercial relational DBMSs employ SQL as their query language.
Alternative query languages have been proposed and implemented, notably the pre-1996 implementation of Ingres QUEL.


== Relational model ==

A relational model organizes data into one or more tables (or "relations") of columns and rows, with a unique key identifying each row. Rows are also called records or tuples. Columns are also called attributes. Generally, each table/relation represents one "entity type" (such as customer or product). The rows represent instances of that type of entity (such as "Lee" or "chair") and the columns represent values attributed to that instance (such as address or price).
For example, each row of a class table corresponds to a class, and a class corresponds to multiple students, so the relationship between the class table and the student table is "one to many"


== Keys ==
Each row in a table has its own unique key. Rows in a table can be linked to rows in other tables by adding a column for the unique key of the linked row (such columns are known as foreign keys). Codd showed that data relationships of arbitrary complexity can be represented by a simple set of concepts.
Part of this processing involves consistently being able to select or modify one and only one row in a table. Therefore, most physical implementations have a unique primary key (PK) for each row in a table. When a new row is written to the table, a new unique value for the primary key is generated; this is the key that the system uses primarily for accessing the table. System performance is optimized for PKs. Other, more natural keys may also be identified and defined as alternate keys (AK). Often several columns are needed to form an AK (this is one reason why a single integer column is usually made the PK). Both PKs and AKs have the ability to uniquely identify a row within a table. Additional technology may be applied to ensure a unique ID across the world, a globally unique identifier, when there are broader system requirements.
The primary keys within a database are used to define the relationships among the tables. When a PK migrates to another table, it becomes a foreign key (FK) in the other table. When each cell can contain only one value and the PK migrates into a regular entity table, this design pattern can represent either a one-to-one or one-to-many relationship. Most relational database designs resolve many-to-many relationships by creating an additional table that contains the PKs from both of the other entity tables –  the relationship becomes an entity; the resolution table is then named appropriately and the two FKs are combined to form a PK. The migration of PKs to other tables is the second major reason why system-assigned integers are used normally as PKs; there is usually neither efficiency nor clarity in migrating a bunch of other types of columns.


=== Relationships ===
Relationships are a logical connection between different tables (entities), established on the basis of interaction among these tables. These relationships can be modelled as an entity-relationship model.


== Transactions ==
In order for a database management system (DBMS) to operate efficiently and accurately, it must use ACID transactions.


== Stored procedures ==
Part of the programming within an RDBMS is accomplished using stored procedures (SPs). Often procedures can be used to greatly reduce the amount of information transferred within and outside of a system. For increased security, the system design may grant access to only the stored procedures and not directly to the tables. Fundamental stored procedures contain the logic needed to insert new and update existing data. More complex procedures may be written to implement additional rules and logic related to processing or selecting the data.


== Terminology ==

The relational database was first defined in June 1970 by Edgar Codd, of IBM's San Jose Research Laboratory. Codd's view of what qualifies as an RDBMS is summarized in Codd's 12 rules. A relational database has become the predominant type of database. Other models besides the relational model include the hierarchical database model and the network model.
The table below summarizes some of the most important relational database terms and the corresponding SQL term:


== Relations or tables ==

In a relational database, a relation is a set of tuples that have the same attributes. A tuple usually represents an object and information about that object. Objects are typically physical objects or concepts. A relation is usually described as a table, which is organized into rows and columns. All the data referenced by an attribute are in the same domain and conform to the same constraints.
The relational model specifies that the tuples of a relation have no specific order and that the tuples, in turn, impose no order on the attributes. Applications access data by specifying queries, which use operations such as select to identify tuples, project to identify attributes, and join to combine relations. Relations can be modified using the insert, delete, and update operators. New tuples can supply explicit values or be derived from a query. Similarly, queries identify tuples for updating or deleting.
Tuples by definition are unique. If the tuple contains a candidate or primary key then obviously it is unique; however, a primary key need not be defined for a row or record to be a tuple. The definition of a tuple requires that it be unique, but does not require a primary key to be defined. Because a tuple is unique, its attributes by definition constitute a superkey.


== Base and derived relations ==
